\section{Introducción}

El Problema del Clique Máximo (\textit{Maximum Clique Problem}, MCP) consiste en encontrar, dado un grafo no dirigido $G=(V,E)$, el subgrafo completo (clique) de mayor cardinalidad, denotada $\omega(G)$ \parencite{pardalos1994maximum}. El problema es NP-duro y, más aún, no admite un algoritmo de aproximación con factor constante en tiempo polinomial salvo que P = NP \parencite{pardalos1994maximum, bomze1999}, lo que motiva el estudio de estrategias de resolución alternativas a la enumeración exacta cuando el tamaño o la densidad de la instancia lo impiden.

En el Informe 1 se abordó el MCP mediante un modelo de Programación Lineal Entera (ILP), recordado aquí brevemente por autocontención: con $x_i\in\{0,1\}$ indicando si el vértice $i$ pertenece al clique, el modelo es $\max \sum_{i \in V} x_i$ sujeto a $x_i + x_j \leq 1\ \forall\ (i,j) \notin E$ (formulación completa y su relajación LP en Informe 1, Sección 3.1). Este modelo fue resuelto de forma exacta con el solver HiGHS \parencite{highs2024} sobre la instancia \textit{keller4} (171 vértices, 9\,435 aristas), obteniendo el óptimo $\omega=11$ en 12--18 segundos. Ese mismo trabajo evidenció, sin embargo, la fragilidad de este enfoque frente a la dificultad de la instancia: la relajación lineal (\textit{Linear Programming}, LP) presentó una brecha de integralidad del 87.1\,\%, y un solver menos sofisticado (GLPK, \textit{GNU Linear Programming Kit}) no logró probar optimalidad dentro de 300 segundos sobre el mismo modelo. Exploraciones adicionales con un Branch \& Bound implementado manualmente sobre una instancia más densa, \textit{C125.9} (125 vértices, densidad 0.90, $\omega=34$), confirmaron esta limitación: tras explorar 232\,953 nodos del árbol de búsqueda durante 300 segundos, el algoritmo encontró el clique óptimo pero \emph{no logró certificar su optimalidad}. Estos antecedentes sitúan el punto de partida de este segundo informe: ¿qué ocurre cuando se abandona la garantía de optimalidad a cambio de un método capaz de escalar a instancias donde el enfoque exacto se degrada?

En la Tarea 1 de la asignatura se diseñaron los componentes fundamentales de una metaheurística de búsqueda local para el MCP: una representación de soluciones como subconjuntos explícitos de vértices que son, por construcción, siempre cliques válidos; una función objetivo trivial $f(S)=|S|$; y un operador de vecindad compuesto por tres movimientos —\textsc{Add}, \textsc{Drop} y \textsc{Swap}— que preserva la factibilidad en cada paso. Este informe construye directamente sobre ese diseño, instanciándolo en una metaheurística concreta: \textbf{Búsqueda Tabú con reinicio por perturbación} (de inspiración en Búsqueda Local Iterada, ILS), y la evalúa computacionalmente para contrastar sus resultados con los del método exacto del Informe 1.

\paragraph{Objetivo general.} Diseñar, implementar y evaluar una metaheurística de búsqueda local para el MCP sobre la vecindad Add/Drop/Swap de la Tarea 1, y contrastar sus resultados —calidad de solución y tiempo de cómputo— con los del método exacto (ILP + Branch \& Bound) del Informe 1.

\paragraph{Objetivos específicos.}
\begin{itemize}
    \item Implementar en Python una Búsqueda Tabú con lista tabú, criterio de aspiración y reinicio por perturbación, sobre la representación y vecindad especificadas en la Tarea 1.
    \item Evaluar el algoritmo mediante corridas independientes sobre \textit{keller4} (la misma instancia del Informe 1) y sobre \textit{C125.9}, una instancia más densa donde el método exacto no logró certificar optimalidad.
    \item Cuantificar la tasa de éxito, el tiempo hasta la mejor solución y la calidad relativa de las soluciones obtenidas respecto al óptimo conocido de cada instancia.
    \item Contrastar sistemáticamente los resultados del enfoque exacto (Informe 1) y del enfoque metaheurístico (este informe), identificando en qué condiciones cada uno resulta preferible.
\end{itemize}

\paragraph{Pregunta de investigación.} ¿En qué medida una metaheurística de búsqueda local relativamente simple, construida sobre una vecindad de tres movimientos que preserva la factibilidad por construcción, logra igualar la calidad de un método exacto para el MCP mientras reduce drásticamente el tiempo de cómputo, y qué se sacrifica —en términos de garantías— al hacerlo?

El resto del documento se organiza como sigue: la Sección~\ref{sec:eda2} presenta el estado del arte de las metaheurísticas para el MCP; la Sección~\ref{sec:met2} describe el método, retomando el diseño de la Tarea 1 y especificando el algoritmo Tabú+ILS; la Sección~\ref{sec:resultados2} presenta los resultados computacionales; la Sección~\ref{sec:contraste} contrasta explícitamente ambos informes; y finalmente se presentan la discusión y las conclusiones.
